% GENERATED FILE. Edit canonical lessons, then run npm run generate:books.
% canonical-source-hash: fnv1a64:1afa4f3e1429e723
% canonical-lessons: MW-C178-ughno, MW-C178-jhakno, MW-C178-binno, MW-C178-bikherno, MW-C178-supno, MW-R178-review-173-178

\chapter{Small Jobs}
\label{ch:mw-small-jobs}
\hlchaptermodality{178}

\begin{chapteropening}
\textbf{By the end of this chapter:} \emph{I can name small jobs around the house.}

Complete the last lesson of chapter 178: I can name small jobs around the house.
\end{chapteropening}

\section[\texorpdfstring{\mw{ऊंघणो} (ū̃ghṇo)}{ऊंघणो (ū̃ghṇo)}]{\mw{ऊंघणो} (ū̃ghṇo) — to doze}

\label{lesson:MW-C178-ughno}



\begin{quote}
Before the new one: say the Marwadi for a patch, then the Marwadi for the end of a sari.
\end{quote}

\subsection*{You'll want to know: \mw{ऊंघणो}}
\textbf{\mw{ऊंघणो}} — \emph{ū̃ghṇo} — \textquotedblleft{}to doze\textquotedblright{}.

\begin{quote}
\mw{दादा} \mw{कुरसी} \mw{माथै} \mw{ऊंघै} \mw{है।}
\end{quote}

\begin{quote}
\emph{dādā kursī māthai ū̃ghai hai.}
\end{quote}

\begin{quote}
Grandfather dozes in the chair.
\end{quote}

Past (he …): \mw{वो} \mw{कुरसी} \mw{माथै} \mw{ऊंघ्यो}, \emph{vo kursī māthai ū̃ghyo}.

\begin{cousinweb}
\textbf{\mw{ऊंघणो}} (\emph{ū̃ghṇo}) comes from a reconstructed \emph{*uṅghati}, \textquotedblleft{}sleeps\textquotedblright{}; CDIAL lists Hindi \emph{ū̃ghnā}, \textquotedblleft{}to be drowsy\textquotedblright{}, and Gujarati \emph{ũghvũ} (CDIAL 1632). CDIAL gives no Marwari form in this entry.
\end{cousinweb}

\begin{grammarlens}[title={What you've built}]
One more word for this chapter.
\end{grammarlens}

\subsection*{Your turn}
\begin{itemize}
\raggedright
  \item \emph{Say it:} \emph{ū̃ghṇo}
  \item \emph{Say it:} \emph{ū̃ghṇo}, once more
  \item \emph{Say it:} \emph{dādā kursī māthai ū̃ghai hai.} — Grandfather dozes in the chair.
  \item \emph{Recall:} say the Marwadi for a patch, then the Marwadi for the end of a sari, then say \emph{ū̃ghṇo} again
\end{itemize}

\begin{tcolorbox}[breakable,colback=teal!4,colframe=teal!35!black,title={Before you move on}]
What does \mw{ऊंघणो} mean? (\textquotedblleft{}To doze\textquotedblright{}.) Say it once more.
\end{tcolorbox}
\section[\texorpdfstring{\mw{झांकणो} (jhā̃kṇo)}{झांकणो (jhā̃kṇo)}]{\mw{झांकणो} (jhā̃kṇo) — to peep}

\label{lesson:MW-C178-jhakno}



\begin{quote}
Before the new one: say the Marwadi for the end of a sari, then the Marwadi for to doze.
\end{quote}

\subsection*{You'll want to know: \mw{झांकणो}}
\textbf{\mw{झांकणो}} — \emph{jhā̃kṇo} — \textquotedblleft{}to peep\textquotedblright{}.

\begin{quote}
\mw{बिल्ली} \mw{खिड़की} \mw{सूं} \mw{झांकै} \mw{है।}
\end{quote}

\begin{quote}
\emph{billī khiṛkī sū̃ jhā̃kai hai.}
\end{quote}

\begin{quote}
The cat peeps in through the window.
\end{quote}

Past (he …): \mw{वो} \mw{खिड़की} \mw{सूं} \mw{झांक्यो}, \emph{vo khiṛkī sū̃ jhā̃kyo}.

\begin{cousinweb}
\textbf{\mw{झांकणो}} (\emph{jhā̃kṇo}) comes from a reconstructed \emph{*jhaṅkh}, \textquotedblleft{}peep\textquotedblright{}; CDIAL lists Old Marwari \emph{jhākaï}, \textquotedblleft{}peeps at\textquotedblright{} (CDIAL 5325).
\end{cousinweb}

\begin{grammarlens}[title={What you've built}]
One more word for this chapter.
\end{grammarlens}

\subsection*{Your turn}
\begin{itemize}
\raggedright
  \item \emph{Say it:} \emph{jhā̃kṇo}
  \item \emph{Say it:} \emph{jhā̃kṇo}, once more
  \item \emph{Say it:} \emph{billī khiṛkī sū̃ jhā̃kai hai.} — The cat peeps in through the window.
  \item \emph{Recall:} say the Marwadi for the end of a sari, then the Marwadi for to doze, then say \emph{jhā̃kṇo} again
\end{itemize}

\begin{tcolorbox}[breakable,colback=teal!4,colframe=teal!35!black,title={Before you move on}]
What does \mw{झांकणो} mean? (\textquotedblleft{}To peep\textquotedblright{}.) Say it once more.
\end{tcolorbox}
\section[\texorpdfstring{\mw{बीणणो} (bīṇṇo)}{बीणणो (bīṇṇo)}]{\mw{बीणणो} (bīṇṇo) — to pick over (grain)}

\label{lesson:MW-C178-binno}



\begin{quote}
Before the new one: say the Marwadi for to doze, then the Marwadi for to peep.
\end{quote}

\subsection*{You'll want to know: \mw{बीणणो}}
\textbf{\mw{बीणणो}} — \emph{bīṇṇo} — \textquotedblleft{}to pick over (grain)\textquotedblright{}.

\begin{quote}
\mw{दादी} \mw{बाजरो} \mw{बीणै} \mw{है।}
\end{quote}

\begin{quote}
\emph{dādī bājro bīṇai hai.}
\end{quote}

\begin{quote}
Grandmother picks over the millet.
\end{quote}

Past (I …): \mw{म्हैं} \mw{बाजरो} \mw{बीण्यो}, \emph{mhaĩ bājro bīṇyo}.

\begin{cousinweb}
\textbf{\mw{बीणणो}} (\emph{bīṇṇo}) is inherited from Sanskrit \emph{vicinoti}, \textquotedblleft{}picks out, chooses\textquotedblright{}; CDIAL lists Hindi \emph{bīnnā}, \textquotedblleft{}to pick up, choose\textquotedblright{} (CDIAL 11686). CDIAL gives no Marwari form in this entry.
\end{cousinweb}

\begin{grammarlens}[title={What you've built}]
One more word for this chapter.
\end{grammarlens}

\subsection*{Your turn}
\begin{itemize}
\raggedright
  \item \emph{Say it:} \emph{bīṇṇo}
  \item \emph{Say it:} \emph{bīṇṇo}, once more
  \item \emph{Say it:} \emph{dādī bājro bīṇai hai.} — Grandmother picks over the millet.
  \item \emph{Recall:} say the Marwadi for to doze, then the Marwadi for to peep, then say \emph{bīṇṇo} again
\end{itemize}

\begin{tcolorbox}[breakable,colback=teal!4,colframe=teal!35!black,title={Before you move on}]
What does \mw{बीणणो} mean? (\textquotedblleft{}To pick over (grain)\textquotedblright{}.) Say it once more.
\end{tcolorbox}
\section[\texorpdfstring{\mw{बिखेरणो} (bikherṇo)}{बिखेरणो (bikherṇo)}]{\mw{बिखेरणो} (bikherṇo) — to scatter}

\label{lesson:MW-C178-bikherno}



\begin{quote}
Before the new one: say the Marwadi for to peep, then the Marwadi for to pick over.
\end{quote}

\subsection*{You'll want to know: \mw{बिखेरणो}}
\textbf{\mw{बिखेरणो}} — \emph{bikherṇo} — \textquotedblleft{}to scatter\textquotedblright{}.

\begin{quote}
\mw{दादी} \mw{बाजरो} \mw{बिखेरै} \mw{है।}
\end{quote}

\begin{quote}
\emph{dādī bājro bikherai hai.}
\end{quote}

\begin{quote}
Grandmother scatters millet.
\end{quote}

Past (I …): \mw{म्हैं} \mw{बाजरो} \mw{बिखेर्यो}, \emph{mhaĩ bājro bikheryo}.

\begin{cousinweb}
\textbf{\mw{बिखेरणो}} (\emph{bikherṇo}) is inherited from a reconstructed \emph{*viṣkirati}, \textquotedblleft{}scatters\textquotedblright{}; CDIAL lists Old Marwari \emph{bakheraï}, \textquotedblleft{}strews\textquotedblright{} (CDIAL 11985).
\end{cousinweb}

\begin{grammarlens}[title={What you've built}]
One more word for this chapter.
\end{grammarlens}

\subsection*{Your turn}
\begin{itemize}
\raggedright
  \item \emph{Say it:} \emph{bikherṇo}
  \item \emph{Say it:} \emph{bikherṇo}, once more
  \item \emph{Say it:} \emph{dādī bājro bikherai hai.} — Grandmother scatters millet.
  \item \emph{Recall:} say the Marwadi for to peep, then the Marwadi for to pick over, then say \emph{bikherṇo} again
\end{itemize}

\begin{tcolorbox}[breakable,colback=teal!4,colframe=teal!35!black,title={Before you move on}]
What does \mw{बिखेरणो} mean? (\textquotedblleft{}To scatter\textquotedblright{}.) Say it once more.
\end{tcolorbox}
\section[\texorpdfstring{\mw{सूंपणो} (sū̃pṇo)}{सूंपणो (sū̃pṇo)}]{\mw{सूंपणो} (sū̃pṇo) — to hand over}

\label{lesson:MW-C178-supno}



\begin{quote}
Before the new one: say the Marwadi for to pick over, then the Marwadi for to scatter.
\end{quote}

\subsection*{You'll want to know: \mw{सूंपणो}}
\textbf{\mw{सूंपणो}} — \emph{sū̃pṇo} — \textquotedblleft{}to hand over\textquotedblright{}.

\begin{quote}
\mw{चाबी} \mw{सूंपो।}
\end{quote}

\begin{quote}
\emph{chābī sū̃po.}
\end{quote}

\begin{quote}
Hand over the key, please.
\end{quote}

Past (I …): \mw{म्हैं} \mw{चाबी} \mw{सूंपी}, \emph{mhaĩ chābī sū̃pī}.

\begin{cousinweb}
\textbf{\mw{सूंपणो}} (\emph{sū̃pṇo}) is inherited from Sanskrit \emph{samarpayati}, \textquotedblleft{}delivers over\textquotedblright{}; CDIAL lists Old Marwari \emph{sopaï}, \textquotedblleft{}hands over\textquotedblright{} (CDIAL 13192).
\end{cousinweb}

\begin{grammarlens}[title={What you've built}]
One more word for this chapter.
\end{grammarlens}

\subsection*{Your turn}
\begin{itemize}
\raggedright
  \item \emph{Say it:} \emph{sū̃pṇo}
  \item \emph{Say it:} \emph{sū̃pṇo}, once more
  \item \emph{Say it:} \emph{chābī sū̃po.} — Hand over the key, please.
  \item \emph{Recall:} say the Marwadi for to pick over, then the Marwadi for to scatter, then say \emph{sū̃pṇo} again
\end{itemize}

\begin{tcolorbox}[breakable,colback=teal!4,colframe=teal!35!black,title={Before you move on}]
What does \mw{सूंपणो} mean? (\textquotedblleft{}To hand over\textquotedblright{}.) Say it once more.
\end{tcolorbox}
\section[(dialogue)]{Looking back: from At the Hearth to Small Jobs}

\label{lesson:MW-R178-review-173-178}



\begin{quote}
Say the Marwadi for to blow on and smoke.
\end{quote}

\subsection*{Your turn}
Say these aloud:

\begin{itemize}
\raggedright
  \item to blow on, smoke, firewood, to put out, to toast
  \item to cover oneself, a head-cloth, a sari, to take off, a plait
  \item a water-carrier, a head-pad, to carry, a small water pot, a spouted water pot
  \item to cover, to tie, a knot, a door chain, to hang up
  \item to tear, a button, to stitch on, a patch, the end of a sari
  \item to doze, to peep, to pick over, to scatter, to hand over
\end{itemize}

\begin{tcolorbox}[breakable,colback=teal!4,colframe=teal!35!black,title={Before you move on}]
To blow on? (\textbf{\mw{फूंकणो}}, \emph{phū̃kṇo}.) Smoke? (\textbf{\mw{धूंवो}}, \emph{dhū̃vo}.) Firewood? (\textbf{\mw{ईंधण}}, \emph{ī̃dhaṇ}.) To put out? (\textbf{\mw{बुझावणो}}, \emph{bujhāvṇo}.) To toast? (\textbf{\mw{सेकणो}}, \emph{sekṇo}.) To cover oneself? (\textbf{\mw{ओढणो}}, \emph{oḍhṇo}.) A head-cloth? (\textbf{\mw{ओढणी}}, \emph{oḍhṇī}.) A sari? (\textbf{\mw{साड़ी}}, \emph{sāṛī}.) To take off? (\textbf{\mw{उतारणो}}, \emph{utārṇo}.) A plait? (\textbf{\mw{चोटी}}, \emph{choṭī}.) A water-carrier? (\textbf{\mw{पणिहारी}}, \emph{paṇihārī}.) A head-pad? (\textbf{\mw{ईंडाणी}}, \emph{ī̃ḍāṇī}.) To carry? (\textbf{\mw{ढोवणो}}, \emph{ḍhovṇo}.) A small water pot? (\textbf{\mw{कळसियो}}, \emph{kaḷsiyo}.) A spouted water pot? (\textbf{\mw{झारी}}, \emph{jhārī}.) To cover? (\textbf{\mw{ढांकणो}}, \emph{ḍhā̃kṇo}.) To tie? (\textbf{\mw{बांधणो}}, \emph{bā̃dhṇo}.) A knot? (\textbf{\mw{गांठ}}, \emph{gā̃ṭh}.) A door chain? (\textbf{\mw{सांकळ}}, \emph{sā̃kaḷ}.) To hang up? (\textbf{\mw{टांगणो}}, \emph{ṭā̃gṇo}.) To tear? (\textbf{\mw{फाड़णो}}, \emph{phāṛṇo}.) A button? (\textbf{\mw{बटण}}, \emph{baṭaṇ}.) To stitch on? (\textbf{\mw{टांकणो}}, \emph{ṭā̃kṇo}.) A patch? (\textbf{\mw{थेगळो}}, \emph{thegḷo}.) The end of a sari? (\textbf{\mw{पल्लो}}, \emph{pallo}.) To doze? (\textbf{\mw{ऊंघणो}}, \emph{ū̃ghṇo}.) To peep? (\textbf{\mw{झांकणो}}, \emph{jhā̃kṇo}.) To pick over? (\textbf{\mw{बीणणो}}, \emph{bīṇṇo}.) To scatter? (\textbf{\mw{बिखेरणो}}, \emph{bikherṇo}.) To hand over? (\textbf{\mw{सूंपणो}}, \emph{sū̃pṇo}.)
\end{tcolorbox}
